\documentclass[10pt,a4paper]{amsart}
\usepackage{amsmath,amssymb,amsthm,graphicx}
% Retained minimal preamble for the Stacks Project chapter more-algebra.tex (More on Algebra).
% The chapter file pinned-original-source.tex carries only \input{preamble} on its first line, so no
% preamble text lives inside the chapter itself.  The Stacks Project uses one single shared preamble for
% every chapter; the snapshot stages it identically as sources/stacks-chow/preamble.tex and
% sources/stacks-etale-cohomology/preamble.tex.  The block below is copied mechanically, line for line,
% from sources/stacks-chow/preamble.tex lines 14-18 (the Xy-pic setup) and 154-217 (the native theorem
% declarations together with the mathematical macros).  Nothing was added, removed or redefined; the
% wrapper class and the amsmath/amssymb/amsthm/graphicx base packages are supplied by the assembler.
% Chapter: more-algebra.tex -- More on Algebra.
% For commutative diagrams we use Xy-pic
\usepackage[all]{xy}

% We use 2cell for 2-commutative diagrams.
\xyoption{2cell}

% Theorem environments.
%
\theoremstyle{plain}
\newtheorem{theorem}[subsection]{Theorem}
\newtheorem{proposition}[subsection]{Proposition}
\newtheorem{lemma}[subsection]{Lemma}

\theoremstyle{definition}
\newtheorem{definition}[subsection]{Definition}
\newtheorem{example}[subsection]{Example}
\newtheorem{exercise}[subsection]{Exercise}
\newtheorem{situation}[subsection]{Situation}

\theoremstyle{remark}
\newtheorem{remark}[subsection]{Remark}
\newtheorem{remarks}[subsection]{Remarks}

\numberwithin{equation}{subsection}

% Macros
%
\def\lim{\mathop{\mathrm{lim}}\nolimits}
\def\colim{\mathop{\mathrm{colim}}\nolimits}
\def\Spec{\mathop{\mathrm{Spec}}}
\def\Hom{\mathop{\mathrm{Hom}}\nolimits}
\def\Ext{\mathop{\mathrm{Ext}}\nolimits}
\def\SheafHom{\mathop{\mathcal{H}\!\mathit{om}}\nolimits}
\def\SheafExt{\mathop{\mathcal{E}\!\mathit{xt}}\nolimits}
\def\Sch{\mathit{Sch}}
\def\Mor{\mathop{\mathrm{Mor}}\nolimits}
\def\Ob{\mathop{\mathrm{Ob}}\nolimits}
\def\Sh{\mathop{\mathit{Sh}}\nolimits}
\def\NL{\mathop{N\!L}\nolimits}
\def\CH{\mathop{\mathrm{CH}}\nolimits}
\def\proetale{{pro\text{-}\acute{e}tale}}
\def\etale{{\acute{e}tale}}
\def\QCoh{\mathit{QCoh}}
\def\Ker{\mathop{\mathrm{Ker}}}
\def\Im{\mathop{\mathrm{Im}}}
\def\Coker{\mathop{\mathrm{Coker}}}
\def\Coim{\mathop{\mathrm{Coim}}}

% Boxtimes
%
\DeclareMathSymbol{\boxtimes}{\mathbin}{AMSa}{"02}

%
% Macros for moduli stacks/spaces
%
\def\QCohstack{\mathcal{QC}\!\mathit{oh}}
\def\Cohstack{\mathcal{C}\!\mathit{oh}}
\def\Spacesstack{\mathcal{S}\!\mathit{paces}}
\def\Quotfunctor{\mathrm{Quot}}
\def\Hilbfunctor{\mathrm{Hilb}}
\def\Curvesstack{\mathcal{C}\!\mathit{urves}}
\def\Polarizedstack{\mathcal{P}\!\mathit{olarized}}
\def\Complexesstack{\mathcal{C}\!\mathit{omplexes}}
% \Pic is the operator that assigns to X its picard group, usage \Pic(X)
% \Picardstack_{X/B} denotes the Picard stack of X over B
% \Picardfunctor_{X/B} denotes the Picard functor of X over B
\def\Pic{\mathop{\mathrm{Pic}}\nolimits}
\def\Picardstack{\mathcal{P}\!\mathit{ic}}
\def\Picardfunctor{\mathrm{Pic}}
\def\Deformationcategory{\mathcal{D}\!\mathit{ef}}
\usepackage{cleveref}
\crefname{originallemma}{original lemma}{original lemmas}
\Crefname{originallemma}{Original lemma}{Original lemmas}
\makeatletter
\newlabel{more-algebra-lemma-quotient-power-series-ring-over-Cohen}{{07NM}{source}{}{sourcepointer-1}{}}
\newlabel{more-algebra-lemma-quotient-power-series-ring-over-Cohen@cref}{{[originallemma][0][]07NM}{[1][0][]source}{}{}{}}
\makeatother
\title{Regular systems of parameters in power series rings over a field and over a Cohen ring}
\author{The Stacks Project Authors}
\date{stacks-project:more-algebra.tex:local-snapshot-20261012; GNU FDL 1.2 or later}
\begin{document}
\maketitle
\setcounter{subsection}{1}
\noindent\textit{Editorial note.} All delivered mathematical text is the author's own, copied from the pinned Stacks Project chapter file more-algebra.tex; no statement or proof step was rewritten, reordered or removed.\par\medskip
\noindent\textit{Editorial note.} The chapter file pinned-original-source.tex carries only \textbackslash{}input\{preamble\} on its first line, so no preamble text lives inside the chapter. The retained file sources/stacks-more-algebra/retained-preamble.tex is a mechanical, line-for-line copy of the project's single shared preamble as staged in the snapshot, namely sources/stacks-chow/preamble.tex lines 14-18 (Xy-pic setup) and 154-217 (theorem declarations and mathematical macros). Nothing was added, removed or redefined, no package was added, and the excerpt serves only as preamble, never as a span coordinate.\par\medskip
\noindent\textit{Editorial note.} Both counted spans start at the \textbackslash{}begin line of their environment, so the statement and the author's proof are delivered whole; surrounding prose of the section is deliberately outside the counted spans and is carried as standard prerequisites instead.\par\medskip
\noindent\textit{Editorial note.} The citation fields (source\_id, source\_version, authors, source\_url, license.name, source\_title) are identical to those of the cards already published from this chapter, so that the batch selector treats them as one source instead of deferring them as a same-source clash.\par\medskip
\noindent\textit{Allowed standard prerequisites (editorial).}\begin{itemize}
\item Regular system of parameters of a Noetherian local ring, and the resulting presentation of the completion. discussed in the section preceding the statement, source lines 9600-9733. Supplies the notion of a regular system of parameters and the structure of a complete regular local ring as a power series ring.
\item Krull intersection or completeness: an isomorphism of complete local rings follows from compatible isomorphisms of all quotients by powers of the maximal ideal. used at proof lines 9756-9758 for part (1). Upgrades the identification of the infinitesimal quotients to the isomorphism K[[y]] to A claimed in (1).
\item A Cohen ring is a complete discrete valuation ring with residue field of characteristic p, with p as part of a regular system of parameters. used in the statement at source lines 9740-9749 and carried by the same argument in (3) and (4). Supplies the second setting in which the same two arguments apply verbatim.
\end{itemize}\medskip
\ifdefined\hypertarget\hypertarget{sourcepointer-1}{}\fi
\noindent\textit{Source reference.} Original lemma 07NM: focused-seeds/AGP/sources/stacks-snapshot-20261011/tags.txt, source lines 9767--9767. The lemma's own canonical tag, quoted so that readers can resolve the label inside the delivered text; it points at the counted statement itself and carries no proof content.\par\medskip
\begin{lemma}
\label{lemma-quotient-power-series-ring-over-Cohen}
Let $K$ be a field and $A = K[[x_1, \ldots, x_n]]$.
Let $\Lambda$ be a Cohen ring and let $B = \Lambda[[x_1, \ldots, x_n]]$.
\begin{enumerate}
\item If $y_1, \ldots, y_n \in A$ is a regular system of parameters
then $K[[y_1, \ldots, y_n]] \to A$ is an isomorphism.
\item If $z_1, \ldots, z_r \in A$ form part of a regular system of
parameters for $A$, then $r \leq n$ and
$A/(z_1, \ldots, z_r) \cong K[[y_1, \ldots, y_{n - r}]]$.
\item If $p, y_1, \ldots, y_n \in B$ is a regular system of parameters
then $\Lambda[[y_1, \ldots, y_n]] \to B$ is an isomorphism.
\item If $p, z_1, \ldots, z_r \in B$ form part of a regular system of
parameters for $B$, then $r \leq n$ and
$B/(z_1, \ldots, z_r) \cong \Lambda[[y_1, \ldots, y_{n - r}]]$.
\end{enumerate}
\end{lemma}

\begin{proof}
Proof of (1). Set $A' = K[[y_1, \ldots, y_n]]$. It is clear that
the map $A' \to A$ induces an isomorphism
$A'/\mathfrak m_{A'}^n \to A/\mathfrak m_A^n$ for all $n \geq 1$.
Since $A$ and $A'$ are both complete we deduce that $A' \to A$ is an
isomorphism.
Proof of (2). Extend $z_1, \ldots, z_r$ to a regular system of parameters
$z_1, \ldots, z_r, y_1, \ldots, y_{n - r}$ of $A$. Consider the map
$A' = K[[z_1, \ldots, z_r, y_1, \ldots, y_{n - r}]] \to A$.
This is an isomorphism by (1). Hence (2) follows as it is clear that
$A'/(z_1, \ldots, z_r) \cong K[[y_1, \ldots, y_{n - r}]]$.
The proofs of (3) and (4) are exactly the same as the proofs of (1) and (2).
\end{proof}

\end{document}
